\begin{proposition}[Bounded measurable functions on finite measure spaces are Lebesgue integrable] \label{proposition:bounded_measurable_functions_on_finite_measure_spaces_are_lebesgue_integrable}
    \TODO{AI generated, verify}
Let $(X,\mathcal{A},\mu)$ be a \CrefAndHyperrefIfExist{definition:measure_on_a_measurable_space_and_measure_space}{measure space} with $\mu(X) < \infty$, and let $f : X \to \overline{\mathbb{R}}$ be a \CrefAndHyperrefIfExist{definition:measurable_function_between_measure_spaces}{measurable function}.

If $f$ is bounded, i.e., there exists $M > 0$ such that $|f(x)| \leq M$ for all $x \in X$, then $f$ is \CrefAndHyperrefIfExist{definition:lebesgue_integral_of_measurable_functions_into_the_real_line_extended_by_infinities}{Lebesgue integrable} with respect to $\mu$, and
$$\int_X |f|\,d\mu \leq M \cdot \mu(X).$$
\end{proposition}

\begin{proof}
    \TODO{AI generated, verify}
Since $|f(x)| \leq M$ for all $x \in X$, we have $|f| \leq M \mathbf{1}_X$. By monotonicity of the Lebesgue integral (if $0 \leq g \leq h$ then $\int_X g\,d\mu \leq \int_X h\,d\mu$),
$$\int_X |f|\,d\mu \leq \int_X M \mathbf{1}_X\,d\mu = M \cdot \mu(X).$$
Since $\mu(X) < \infty$ and $M > 0$, the right-hand side is finite. Hence $\int_X |f|\,d\mu < \infty$, which means $f$ is Lebesgue integrable by definition.
\end{proof}
